\appendix
\section{Appendix}\label{sec:appendix}

\subsection{Temporal predicates}\label{sec:app-temporal}
The order and containment predicates of \cref{sec:chaining} compare interval endpoints under a slack
of one quarter of the shorter window; for windows $W_X=[s_X,e_X]$,
\begin{equation}
\begin{gathered}
\delta(A,B)=\tfrac14
  \min(e_A-s_A,e_B-s_B),\\
\rho_{\mathrm{before}}(W_A,W_B)
=
\begin{cases}
\mathsf{S}, & e_A\leq s_B+\delta(A,B),\\
\mathsf{C}, & e_B\leq s_A+\delta(A,B),\\
\mathsf{U}, & \text{otherwise}.
\end{cases}
\end{gathered}
\label{eq:temporal-before}
\end{equation}
The quarter-window rule keeps the first two cases mutually exclusive for positive durations;
containment uses the same $\delta$.

\subsection{Withdrawal noise-level schedule}\label{sec:app-schedule}
\Cref{sec:gen-withdrawal} reports the withdrawal point by noise level $\sigma_{\mathrm{rel}}$, not
step index $i^{*}$. For $S$ steps at flow shift $\kappa$ over $N$ training timesteps, the UniPC
schedule assigns
\begin{equation}
\tilde{s}_i = \frac{(S-i)\,(1 - 1/N)}{S},
\qquad
\sigma_i = \frac{\kappa\,\tilde{s}_i}{1 + (\kappa-1)\,\tilde{s}_i},
\qquad
\sigma_{\mathrm{rel}} = \sigma_{i^{*}} ,
\label{eq:flow-sigma}
\end{equation}
with $\tilde{s}_i$ the normalized flow-matching timestep at step $i$ and $\sigma_i$ its noise level.

\subsection{Trajectory-agreement metrics}\label{sec:app-retention}
\Cref{sec:gen-release} scores how closely a generated clip's tracked path follows the simulation's.
With $p^{\mathrm{gen}}_t,p^{\mathrm{ctrl}}_t$ the per-frame tracked positions ($t\in\{0,\ldots,T-1\}$,
undefined where the tracker finds nothing) and $V$ the frames where both are defined,
\emph{retention} over a window $W$ is the per-axis Pearson correlation over $W\cap V$ ($R_x(W)$,
$R_y(W)$), reported only where $|W\cap V|$ meets a minimum count; $R$ near $1$ means the object moved
as simulated.

The windows split at the first ground-contact frame, off the control's vertical path: with
$y^{\mathrm{ctrl}}$ downward, $t_{\mathrm{apex}}=\arg\min_{t\in V} y^{\mathrm{ctrl}}_t$, and
$y_{\min},y_{\max}$ over the located frames,
$t_c=\min\{t\geq t_{\mathrm{apex}}: y^{\mathrm{ctrl}}_t \geq y_{\max}-\tfrac{1}{10}(y_{\max}-y_{\min})\}$,
giving \emph{airborne} $\{t<t_c\}$ and \emph{after-contact} $\{t\geq t_c\}$. Two prefixed failure
modes accompany retention: a \emph{whole-clip disagreement} (full-clip horizontal $<0.7$) and a
\emph{late-window disagreement} (full-clip horizontal $<0.7$ while airborne horizontal $>0.9$).

\subsection{Critic dataflow and procedure}\label{sec:app-critic-alg}
\Cref{fig:criticpipe} shows the dataflow of \cref{sec:critic} and \cref{alg:critic} the full one-clip
procedure, each step tagged learned or deterministic with its defining equation; both summarise
\cref{sec:critic} and add no new step.

\begin{figure}[H]
  \centering
  % =====================================================================
% fig_critic_pipeline.tex -- dataflow of the auditable multi-tool critic.
%
% \input-able into main.tex, which loads adobe-techreport.sty (hence the
% adobered / adobeink / adobegray colours and the Source font family).
%
% REQUIRED IN THE MAIN PREAMBLE:
%     \usepackage{tikz}
%     \usetikzlibrary{arrows.meta,positioning,calc}
% Both lines are required, and the libraries MUST be loaded in the preamble --
% see the guard below for why a figure file cannot load them for itself.
% Every style is declared inside the tikzpicture's own option list, so all
% definitions are local to this picture: \input-ing this file twice, or
% alongside its companion figure, cannot clash.
%
% CONTENT follows the section "The Auditable Multi-Tool Critic": the overview,
% the reasoner, the perception specialists, the evidence-only tool contract,
% validation of a plan before any tool runs, and the three-valued semantics
% with a deterministic roll-up. The figure is descriptive of the method only:
% no accuracy, no coverage, no comparison between configurations, and no
% quantity of any kind appears in it.
%
% Measured width when \input into the report is 409.6pt against a 469.8pt
% \textwidth, so it fits unscaled -- do not add scale= or \resizebox. The build
% script re-checks the width, and re-checks that the width is identical on a
% second inclusion:
%     bash /Users/wenzhuox/diffphy_psc/techreport/figs/build_critic_figs.sh
% =====================================================================
% The three TikZ libraries must be loaded in the PREAMBLE, not here: a figure
% environment forms a group, and TikZ defines a library's macros locally, so a
% \usetikzlibrary issued inside a figure does not survive to the next one. This
% guard turns a missing preamble line into a clear message instead of TikZ's
% confusing "You need to say \usetikzlibrary{calc}" further down.
\makeatletter
\@ifundefined{tikz@library@calc@loaded}{%
  \PackageError{fig_critic_pipeline}{%
    The preamble must contain\MessageBreak
    \string\usetikzlibrary{arrows.meta,positioning,calc}%
  }{Add that line after \string\usepackage{tikz}.}}{}
\makeatother

\begingroup
\providecommand{\cfsub}[1]{{\scriptsize\color{adobegray}#1}}

\begin{tikzpicture}[
    x=1cm, y=1cm,
    % --- shared box shape -------------------------------------------------
    cfbox/.style   = {rounded corners=2pt, line width=0.7pt, align=center,
                      inner sep=5pt, font=\sffamily\footnotesize,
                      text=adobeink},
    % --- LEARNED component: light gray fill --------------------------------
    cflearn/.style = {cfbox, draw=adobeink, fill=adobegray!12},
    % --- DETERMINISTIC component: white with an ink outline ----------------
    cfdet/.style   = {cfbox, draw=adobeink, fill=white},
    % --- accent 1: the deterministic check that precedes every tool call ---
    cfgate/.style  = {cfbox, draw=adobered, fill=white},
    % --- the critic's verdict (drawn like any deterministic box) -----------
    cfout/.style   = {cfbox, draw=adobeink, fill=white},
    % --- what enters the critic --------------------------------------------
    cfdata/.style  = {cfbox, draw=adobegray, fill=white, inner sep=4pt,
                      dash pattern=on 2pt off 1.6pt},
    cfline/.style  = {draw=adobeink, line width=0.7pt},
    cfarr/.style   = {cfline, -{Straight Barb[length=3.6pt,width=3.8pt]}},
    cfnote/.style  = {font=\sffamily\scriptsize, text=adobegray,
                      inner sep=0pt, align=left},
    cfsw/.style    = {rounded corners=1.2pt, line width=0.7pt, inner sep=0pt,
                      minimum width=11pt, minimum height=7.5pt},
  ]

  % Full-width rows sit on the x=0 axis; the two-column rows are anchored to
  % the full-width box above them, so every left and right edge lines up.

  % ---------------- reasoner (learned), fed by the prompt and the video ----
  % This row plus its two input wings is the widest thing in the picture, so
  % its width is what keeps the whole figure inside the report's text width.
  \node[cflearn, text width=9.4cm, anchor=north] (reasoner) at (0,0)
    {\textbf{Reasoner}\\[1.5pt]
     \cfsub{breaks the prompt into small checkable claims and writes down,\\
     for each claim, the measurements that would settle it ---\\
     without looking at the video yet}};

  \node[cfdata, text width=1.25cm, anchor=east]
    (prompt) at ($(reasoner.west)+(-0.62,0.20)$) {\textbf{Prompt}};
  \node[cfdata, text width=1.25cm, anchor=west]
    (video)  at ($(reasoner.east)+( 0.62,0.20)$) {\textbf{Video}};

  % ---------------- the check that precedes execution (accented) -----------
  \node[cfgate, text width=13.6cm, anchor=north]
    (gate) at ($(0,0)!1!(reasoner.south)+(0,-0.60)$)
    {{\color{adobered}\textbf{Every plan is checked before any tool runs}}\\[1.5pt]
     \cfsub{each line is parsed and type-checked; a line that cannot be used
     falls back to a general judgment,\\ and that fallback is counted in the
     plan's cost rather than treated as free}};

  % ---------------- execution: two ways of looking at the video -----------
  \coordinate (fanA) at ($(gate.south)+(0,-0.40)$);
  \node[cflearn, text width=5.6cm, anchor=north west]
    (verifier) at ($(gate.south west)+(0,-0.80)$)
    {\textbf{General verifier}\\[1.5pt]
     \cfsub{reads the frames densely and decides\\ whether a thing is there and
     whether\\ an event happens}};
  \node[cflearn, text width=7.0cm, anchor=north east]
    (specialists) at ($(gate.south east)+(0,-0.80)$)
    {\textbf{Perception specialists}\\[1.5pt]
     \cfsub{one frozen model per measurement: detection,\\ counting, point
     tracking, the time window of\\ an event, segmentation masks with depth}};

  % ---------------- evidence-only contract (deterministic) ---------------
  \coordinate (mergeA) at ($(verifier.south)+(0,-0.50)$);
  \node[cfdet, text width=13.6cm, anchor=north west]
    (evidence) at ($(verifier.north west |- mergeA)+(0,-0.40)$)
    {\textbf{Evidence only}\\[1.5pt]
     \cfsub{a call returns what it measured --- a box, a count, a track, a time
     window, a depth --- tied to the\\ prompt fragment that asked for it, and
     never whether the prompt was satisfied}};

  % ---------------- deterministic composition ----------------------------
  \node[cfdet, text width=13.6cm, anchor=north west]
    (compose) at ($(evidence.south west)+(0,-0.40)$)
    {\textbf{Deterministic composition}\\[1.5pt]
     \cfsub{fixed rules combine the parts of a claim, then the claims into the
     clip: anything contradicted makes\\ its parent contradicted, everything
     supported makes it supported, and anything else is left unresolved}};

  % ---------------- outputs ----------------------------------------------
  \coordinate (fanB) at ($(compose.south)+(0,-0.40)$);
  \node[cfout, text width=5.6cm, anchor=north west]
    (verdict) at ($(compose.south west)+(0,-0.80)$)
    {{\textbf{Verdict}}\\[1.5pt]
     \cfsub{each claim supported, contradicted or\\ unknown; the whole clip
     plausible,\\ implausible, or abstained on}};
  \node[cfdet, text width=7.0cm, anchor=north east]
    (trace) at ($(compose.south east)+(0,-0.80)$)
    {\textbf{Audit trace}\\[1.5pt]
     \cfsub{the claims checked, the evidence each call\\ returned, which claims
     fell back to a general\\ judgment, and the rule that combined them}};

  % ---------------- edges ------------------------------------------------
  \draw[cfarr] (prompt.east) -- (reasoner.west);
  \draw[cfarr] (video.west)  -- (reasoner.east);
  \draw[cfarr] (reasoner.south) -- (reasoner.south |- gate.north);

  \draw[cfline] (gate.south) -- (fanA);
  \draw[cfline] (verifier.north |- fanA) -- (specialists.north |- fanA);
  \draw[cfarr]  (verifier.north |- fanA) -- (verifier.north);
  \draw[cfarr]  (specialists.north |- fanA) -- (specialists.north);

  % both branches drop to one horizontal line, which then feeds the evidence box
  \draw[cfline] (verifier.south)    -- (verifier.south    |- mergeA);
  \draw[cfline] (specialists.south) -- (specialists.south |- mergeA);
  \draw[cfline] (verifier.south |- mergeA) -- (specialists.south |- mergeA);
  \draw[cfarr]  (evidence.north |- mergeA) -- (evidence.north);

  \draw[cfarr] (evidence.south) -- (evidence.south |- compose.north);

  \draw[cfline] (compose.south) -- (fanB);
  \draw[cfline] (verdict.north |- fanB) -- (trace.north |- fanB);
  \draw[cfarr]  (verdict.north |- fanB) -- (verdict.north);
  \draw[cfarr]  (trace.north   |- fanB) -- (trace.north);

  % ---------------- inline legend ----------------------------------------
  % Deliberately terse: this row is the widest thing in the picture, and the
  % whole picture must stay inside the report's text width.
  \node[cfsw, draw=adobeink, fill=adobegray!12, anchor=north west]
    (swA) at ($(verdict.south west)+(0,-0.50)$) {};
  \node[cfnote, right=3pt of swA] (lbA) {learned};

  \node[cfsw, draw=adobeink, fill=white, right=9pt of lbA] (swB) {};
  \node[cfnote, right=3pt of swB] (lbB) {deterministic};

  \node[cfsw, draw=adobered, fill=white, right=9pt of lbB] (swC) {};
  \node[cfnote, right=3pt of swC] (lbC) {deterministic, before any tool};

  \node[cfsw, draw=adobegray, fill=white, dash pattern=on 2pt off 1.6pt,
        right=9pt of lbC] (swE) {};
  \node[cfnote, right=3pt of swE] (lbE) {input};

\end{tikzpicture}
\endgroup

  \caption{\PhyReAct{} dataflow. The text-only planner writes claims and a surface
  plan; the compiler expands and validates it before any video is read. SAM~3
  grounds the subject and supplies the tracks the typed operations measure over, so
  specialists return numbers while the video-aware verifier returns tagged learned
  states. Composition is deterministic; the one learned step after it, the fallback
  verifier, may resolve a claim the dense pass left supported. Shaded boxes are
  learned, outlined boxes deterministic.}
  \label{fig:criticpipe}
\end{figure}

\begin{algorithm}[H]
\caption{VeriPhy critic: bounded learned planning/verification embedded in deterministic compilation, scoping, and roll-up (executable summary of \cref{sec:critic}).}
\label{alg:critic}
\begin{algorithmic}[1]
\Require clip $\mathcal{V}=\{(I_t,\tau_t)\}_{t\in\mathcal{T}_{\mathcal V}}$, prompt $p$
\Ensure verdict $Y\in\{\textsc{Plausible},\textsc{Implausible},\textsc{Abstain}\}$ (or $\bot_{\mathrm{infra}}$)
\Procedure{Critic}{$\mathcal{V},p$}
  \State $(\mathcal{C},\widetilde{\Pi})\gets\Pi_\theta(p)$, \; $\mathcal{C}=\{c_i=(s_i,\phi_i)\}_{i=1}^{N_p}$, $s_i\sqsubseteq p$ \Comment{\cref{eq:claims-and-plan}; text-only, video unseen \emph{(learned)}}
  \State $G_\Pi=(\mathcal{A},\mathcal{Q},\mathcal{E},b)\gets\textsc{Compile}(\widetilde{\Pi})$, \; $a_j=(o_j,\alpha_j,\omega_j,r_j)$ \Comment{typed DAG, before any frame is read; \cref{eq:plan-dag} \emph{(det.)}}
  \If{$\neg\operatorname{Valid}(G_\Pi)$} \Comment{$\operatorname{Valid}=\operatorname{Parse}\land\operatorname{Refs}\land\operatorname{Typed}\land\operatorname{Acyclic}\land[\forall i\,\exists q\!:b(q){=}i]$; \cref{eq:plan-validity} \emph{(det.)}}
     \State $G_\Pi\gets\textsc{Fallback}(G_\Pi)$; \; \textsc{Charge}\,\&\,\textsc{Report} \Comment{invalid line$\to$general verifier, invalid term$\to$recorded sub-check; never guessed}
  \EndIf
  \For{$a_j\in\mathcal{A}$ in topological order of $G_\Pi$} \Comment{video first touched here}
     \State $\Omega_j\gets g_j(\{R_k:(a_k,a_j)\in\mathcal{E}\},\mathcal{T}_{\mathcal V})$; \; $R_j\gets F_{o_j}(\mathcal{V}|_{\Omega_j},\alpha_j)$ \Comment{obs.\ gate scope \emph{(det.)}; frozen backend, SAM~3 grounds tracks; \cref{eq:scoped-execution}}
  \EndFor
  \For{$q\in\mathcal{Q}$}
     \State $z_q\gets\psi_q(\{\mu_j:a_j\in\operatorname{Pa}(q)\})\in\mathbb{V}=\{\mathsf{S},\mathsf{C},\mathsf{U}\}$; \; negate via $\neg_3$ \Comment{$\psi_q$ det.\ for count/timing/track/transcription, learned asymmetric gate ($\mathsf{C}$ only on positive contradiction) for semantic; $\mathsf{U}$ if abstained, $\bot_{\mathrm{infra}}$ on infra error; \cref{sec:semantics}}
  \EndFor
  \For{$i\gets1$ to $N_p$}
     \State $K_i\gets\{q:b(q)=i\}$; \; $y_i\gets\mathsf{C}$ if $\exists q\!\in\!K_i\!:z_q{=}\mathsf{C}$, \; $\mathsf{S}$ if $\forall q\!\in\!K_i\!:z_q{=}\mathsf{S}$, \; else $\mathsf{U}$ \Comment{claim roll-up; \cref{eq:claim-clip-rollup} \emph{(det.)}}
  \EndFor
  \ForAll{$i$ with $y_i=\mathsf{U}$} \Comment{single post-composition step \emph{(learned)}}
     \State $y_i\gets\textsc{FallbackVerify}(c_i,\{R_j\})$
  \EndFor
  \State $Y\gets\bot_{\mathrm{infra}}$ on infra error; \; \textsc{Implausible} if $\exists i\!:y_i{=}\mathsf{C}$; \; \textsc{Plausible} if $N_p{\ge}1\land\forall i\,y_i{=}\mathsf{S}$; \; else \textsc{Abstain} \Comment{clip roll-up; \cref{eq:claim-clip-rollup} \emph{(det.)}}
  \State \Return $Y$
\EndProcedure
\end{algorithmic}
\end{algorithm}

\begin{figure}[tp]
  \centering
  % =====================================================================
% fig_plan_graph.tex -- one prompt's verification plan as a graph:
% claims -> typed checks -> the measurements each check consumes.
% \input into main.tex, which loads adobe-techreport.sty (adobered /
% adobeink / adobegray) and tikz with arrows.meta, positioning, calc.
% =====================================================================
{%
\ifdefined\tikzlibraryarrowsmetaloaded\else
  \expandafter\ifx\csname tikz@library@arrows.meta@loaded\endcsname\relax
    \PackageError{fig_plan_graph}{The preamble must contain
      \string\usetikzlibrary{arrows.meta,positioning,calc}}{}%
  \fi
\fi
%
\footnotesize
\begin{tikzpicture}[
  font=\footnotesize,
  >={Stealth[length=4pt,width=3pt]},
  node distance=0pt,
  % a claim taken from the prompt
  pgclaim/.style   = {draw=adobeink, fill=adobegray!14, rounded corners=1.6pt,
                      line width=0.6pt, align=left, inner sep=4pt,
                      text width=3.32cm, minimum height=0.86cm},
  % a typed check over measurements
  pgcheck/.style   = {draw=adobered, fill=white, rounded corners=1.6pt,
                      line width=0.6pt, align=center, inner sep=3.2pt,
                      text width=2.32cm},
  % a measurement handed to a tool
  pgmeas/.style    = {draw=adobeink, fill=white, rounded corners=1.6pt,
                      line width=0.6pt, align=left, inner sep=3.2pt,
                      text width=4.15cm},
  % a measurement answered by the general verifier rather than a specialist
  pggen/.style     = {pgmeas, fill=adobegray!8, draw=adobegray},
  pgedge/.style    = {draw=adobegray, line width=0.55pt, ->},
  pglab/.style     = {font=\scriptsize, text=adobegray, inner sep=1.2pt},
]

% ---- the plan, as the planner wrote it -------------------------------
\node[draw=adobegray, fill=adobegray!6, line width=0.5pt, rounded corners=1.6pt,
      inner sep=5pt, text width=13.05cm, align=left, font=\ttfamily\scriptsize]
  (code) at (0,0)
  {check(c0, judge("a dolphin leaps out of the water"))\\
   check(c1, before(window("the dolphin leaps"), window("the dolphin dives back under")))\\
   check(c2, judge("two salmon emerge at the same moment") and eq(count("salmon"), 2))\\
   check(c3, judge("the salmon are behind the fishing boat") and rel("salmon||fishing boat", "behind"))};

% ---- row 1: a single occurrence check --------------------------------
\node[pgclaim, anchor=north west] (c0) at ($(code.south west)+(0.10,-0.62)$)
  {\textbf{c0}\\ the dolphin leaps};
\node[pgcheck, anchor=west] (p0) at ($(c0.east)+(1.05,0)$) {event occurs};
\node[pggen, anchor=west]  (m0) at ($(p0.east)+(1.05,0)$)
  {general verifier};
\draw[pgedge] (c0) -- (p0);
\draw[pgedge] (p0) -- (m0);

% ---- row 2: ordering, two measurements into one check ----------------
\node[pgclaim, anchor=north west] (c1) at ($(c0.south west)+(0,-0.55)$)
  {\textbf{c1}\\ it dives back \emph{after} the leap};
\node[pgcheck, anchor=west] (p1) at ($(c1.east)+(1.05,0)$) {order in time\\ (before, with slack)};
\node[pgmeas, anchor=south west] (m1a) at ($(p1.east)+(1.05,0.06)$)
  {event timing (leap)};
\node[pgmeas, anchor=north west] (m1b) at ($(p1.east)+(1.05,-0.06)$)
  {event timing (dive)};
\draw[pgedge] (c1) -- (p1);
\draw[pgedge] (p1.east) to[out=25,in=180] (m1a.west);
\draw[pgedge] (p1.east) to[out=-25,in=180] (m1b.west);

% ---- row 3: a claim carrying two checks ------------------------------
\node[pgclaim, anchor=north west] (c2) at ($(c1.south west)+(0,-1.15)$)
  {\textbf{c2}\\ two salmon emerge together};
\node[pgcheck, anchor=south west] (p2a) at ($(c2.east)+(1.05,0.42)$) {event occurs};
\node[pgcheck, anchor=north west] (p2b) at ($(c2.east)+(1.05,-0.42)$) {count is 2};
\node[pggen, anchor=west] (m2a) at ($(p2a.east)+(1.05,0)$)
  {general verifier};
\node[pgmeas, anchor=west] (m2b) at ($(p2b.east)+(1.05,0)$)
  {counter};
\draw[pgedge] (c2.east) to[out=25,in=180] node[pglab,pos=0.55,above]{and} (p2a.west);
\draw[pgedge] (c2.east) to[out=-25,in=180] node[pglab,pos=0.55,below]{and} (p2b.west);
\draw[pgedge] (p2a) -- (m2a);
\draw[pgedge] (p2b) -- (m2b);

% ---- row 4: a claim whose second check needs a specialist ------------
\node[pgclaim, anchor=north west] (c3) at ($(c2.south west)+(0,-1.45)$)
  {\textbf{c3}\\ the salmon are behind the boat};
\node[pgcheck, anchor=south west] (p3a) at ($(c3.east)+(1.05,0.42)$) {event occurs};
\node[pgcheck, anchor=north west] (p3b) at ($(c3.east)+(1.05,-0.42)$) {behind\\ (depth order)};
\node[pggen, anchor=west] (m3a) at ($(p3a.east)+(1.05,0)$)
  {general verifier};
\node[pgmeas, anchor=west] (m3b) at ($(p3b.east)+(1.05,0)$)
  {outlines and depth};
\draw[pgedge] (c3.east) to[out=25,in=180] node[pglab,pos=0.55,above]{and} (p3a.west);
\draw[pgedge] (c3.east) to[out=-25,in=180] node[pglab,pos=0.55,below]{and} (p3b.west);
\draw[pgedge] (p3a) -- (m3a);
\draw[pgedge] (p3b) -- (m3b);

% ---- column headings --------------------------------------------------
\node[pglab, anchor=south] at ($(c0.north)+(0,0.10)$) {claims from the prompt};
\node[pglab, anchor=south] at ($(p0.north)+(0,0.10)$) {typed checks};
\node[pglab, anchor=south] at ($(m0.north)+(0,0.10)$) {measurements};

% ---- key -------------------------------------------------------------
\node[pgclaim, anchor=north west, text width=0.30cm, minimum height=0.30cm,
      inner sep=0pt] (kA) at ($(c3.south west)+(0,-1.05)$) {};
\node[pglab, anchor=west] (kAl) at ($(kA.east)+(0.10,0)$) {a claim};
\node[pgcheck, anchor=west, text width=0.30cm, minimum height=0.30cm,
      inner sep=0pt] (kB) at ($(kAl.east)+(0.55,0)$) {};
\node[pglab, anchor=west] (kBl) at ($(kB.east)+(0.10,0)$) {a typed check};
\node[pggen, anchor=west, text width=0.30cm, minimum height=0.30cm,
      inner sep=0pt] (kC) at ($(kBl.east)+(0.55,0)$) {};
\node[pglab, anchor=west] (kCl) at ($(kC.east)+(0.10,0)$) {answered by the general verifier};
\node[pgmeas, anchor=west, text width=0.30cm, minimum height=0.30cm,
      inner sep=0pt] (kD) at ($(kCl.east)+(0.55,0)$) {};
\node[pglab, anchor=west] at ($(kD.east)+(0.10,0)$) {measured by a specialist};

\end{tikzpicture}%
}

  \caption{A surface plan compiled into claims, typed checks, and evidence-producing
  actions (from \cref{sec:plan-lang}): a numeral adds a count check beside a semantic
  assertion; an ordering consumes two event windows after occurrence gating; a depth
  relation consumes localization, masks, and depth from the compiler's dependency
  closure. Illustrative of the planner's vocabulary, not an added learned verdict.}
  \label{fig:plangraph}
\end{figure}

% =====================================================================
\subsection{Benchmark details}\label{sec:app-benchmark}
Field coverage is in \cref{tab:flawfields}, a full record in \cref{fig:benchanatomy}, the set's
derivation in \cref{fig:benchcomp}, and the core's distributions in \cref{fig:benchdist}.

By machine-assigned type, the core's $306$ flaws (median two per clip) split as actions absent or
wrong $139$; wrong counts $36$; wrong or missing objects $30$; garbled text $25$; wrong spatial
arrangement $19$; wrong ordering $17$; wrong attributes $16$; camera and style $10$; physical motion
$7$; residual $7$---counting, text, and arrangement together $80$, so this is not a physics benchmark
in the narrow sense. The $150$ clips run a median $8.0$\,s ($5.0$--$10.3$), mostly $24$\,fps
($118/150$), under prompts of median $42$ words (longest $188$), many naming a shot list, count,
ordering, on-screen string, and sound at once. Sixty-six carry more than one flaw kind. Of the $237$
time-spanned flaws the flagged window is bimodal---$57$ span a quarter of the clip or less, $95$ at
least nine-tenths (\cref{fig:benchdist})---so a detector needs both a localized and a whole-clip
check.

\begin{table}[t]
\centering\small
\setlength{\tabcolsep}{6pt}\renewcommand{\arraystretch}{1.18}
\begin{tabular}{@{}l l c@{}}
\toprule
\textbf{Field} & \textbf{Note} & \textbf{Coverage} \\
\midrule
violated fragment & median $8$ words & 306/306 \\
rationale & median $25$ words & 306/306 \\
severity & $1$--$5$, centred at $3$--$4$ & 306/306 \\
confidence & $1$--$5$, $239$ marked $4$--$5$ & 306/306 \\
time span & & 237/306 \\
box track & & 208/306 \\
\bottomrule
\end{tabular}
\caption{Coverage of a flaw record's six fields (defined in \cref{sec:bench-spec}) across the
$150$-clip core: the first four always present, timing and boxes only when the annotator
supplies them.}
\label{tab:flawfields}
\end{table}

\begin{figure}[t]
  \centering
  \includegraphics[width=\linewidth]{fig_benchmark_anatomy.pdf}
  \caption{One human flaw record in full: the quoted prompt fragment, the interval
  of visible failure ($1.62$\,s of an $8$\,s clip), boxes on the objects at the
  interval's start/middle/end, and a severity, confidence, and rationale. No critic
  output appears.}
  \label{fig:benchanatomy}
\end{figure}

\begin{figure}[t]
  \centering
  \includegraphics[width=\linewidth]{fig_benchmark_composition.pdf}
  \caption{Deriving the evaluation set. Only the first step removes clips; the
  second removes audible-but-not-visible flaw \emph{records}, leaving the clip
  count unchanged. The core is a fixed random sample balanced by severity, flaw
  count, and length; a further $711$ annotated clips, disjoint from both frozen
  sets, are kept for training.}
  \label{fig:benchcomp}
\end{figure}

\begin{figure}[t]
  \centering
  \includegraphics[width=\linewidth]{fig_benchmark_distributions.pdf}
  \caption{Two core properties that shape a detector. \textbf{Left:} most clips are
  flagged in only one or two places, so many complaints per clip are mostly ones no
  annotator marked. \textbf{Right:} flagged windows are bimodal---some failures
  confined to a small part of the clip, a comparable group spanning nearly all of
  it---so both a localized and a whole-clip check are needed.}
  \label{fig:benchdist}
\end{figure}

\subsection{Published-evaluator comparison, in full}\label{sec:app-external}
\Cref{tab:external} is the full form of \cref{fig:external}, adding each evaluator's output and the
modular video-QA row (listed for context, not evaluated under this protocol, so no figure point).

% GENERATED by analysis/emit_tex.py — do not edit by hand.
% Every number is recomputed from per-item rows; see the analysis scripts.

\begin{table}[t]
\centering\small
\setlength{\tabcolsep}{5pt}\renewcommand{\arraystretch}{1.18}
\begin{tabular}{@{}p{4.5cm} p{3.3cm} c c c@{}}
\toprule
\textbf{Evaluator} & \textbf{Output it produces} & \textbf{Found} & \textbf{Recall}
 & \textbf{Calls} \\
\midrule
This work & localized allegation & \textbf{228} & \textbf{75.0\%} & 14.1 \\
\addlinespace
\multicolumn{5}{@{}l}{\textit{Comparable on localized recall}} \\
\quad Question decomposition~\cite{tifa2023} & per-question answers & 164 & 53.9\% & 21.9 \\
\quad Davidsonian scene graph~\cite{dsg2024} & per-question answers, dependency-gated & 161 & 53.0\% & 52.1 \\
\quad Modular video QA~\cite{proviq2023,morevqa2024} & per-question answers & --- & --- & --- \\
\addlinespace
\multicolumn{5}{@{}l}{\textit{Recall reported, but not a stable quantity}} \\
\quad Single-pass VLM, same model~\cite{qwen3vl2025} & free-text flaw list & 222 & 73.0\% & 1.0 \\
\bottomrule
\end{tabular}
\caption{Localized flaw recall over the same 150 clips and the same $304$
flaws, judged by the protocol of \cref{sec:matching}. Cost counts model and tool calls per clip and
excludes the matching protocol, which is the scorer rather than the evaluator. Question
decomposition is given the same claim decomposition \PhyReAct{} used, so its recall cannot differ
because it checked different things. The single-pass row is the same served model as \PhyReAct{}'s
own, prompted once with every frame; its recall is set apart because rephrasing the prompt changes
the answer. A dash marks an evaluator nobody has run under this protocol.}
\label{tab:external}
\end{table}


\subsection{Physion-Eval comparison, in full}\label{sec:app-physion}
\Cref{tab:physion} is the full form of \cref{fig:physionlb}, adding this work's matched-report and
anomaly-phrased rates and the benchmark's $J$ statistics---defined against withheld real-world
reference videos, hence dashed for every evaluator on the release.

% Numbers recomputed from per-clip rows by video_agentic_critic/analysis/physion_section5_numbers.py
% (this work) and transcribed from Physion-Eval's Figure 4 (the benchmark's critics).

\begin{table}[t]
\centering\small
\setlength{\tabcolsep}{6pt}\renewcommand{\arraystretch}{1.18}

\begin{tabular}{@{}p{6.2cm} c c@{}}
\multicolumn{3}{@{}l}{\textbf{Panel A: this work, on the interim stream ($601$ release clips)}} \\
\toprule
\textbf{Measure} & \textbf{All clips ($n=601$)} & \textbf{Reported-glitch clips ($n=543$)} \\
\midrule
Evidence-conditioned flag rate & $26.8\%$ & $27.8\%$ \\
Anomaly-phrased finding rate & $15.5\%$ & $15.3\%$ \\
Matched-report clip rate & --- & $18.4\%$ \\
\bottomrule
\end{tabular}

\vspace{1.2em}

\begin{tabular}{@{}p{6.2cm} c c c@{}}
\multicolumn{4}{@{}l}{\textbf{Panel B: benchmark-reported results, on the benchmark's own evaluation set}} \\
\toprule
\textbf{Evaluator} & \textbf{$J$ exo} & \textbf{$J$ ego} & \textbf{Flag rate} \\
\midrule
Untrained humans (reference) & $30.7\%$ & $56.0\%$ & $52.7\%$ \\
Gemini 3.0 Pro & $11.3\%$ & $5.5\%$ & $16.9\%$ \\
Qwen3-VL-32B & $3.2\%$ & $1.0\%$ & $14.4\%$ \\
Qwen3-VL-8B & $2.7\%$ & $0.8\%$ & $12.7\%$ \\
Claude 4.5 Opus & $5.7\%$ & $2.3\%$ & $10.7\%$ \\
GPT-5.2 & $4.0\%$ & $3.7\%$ & $10.5\%$ \\
Gemini 2.5 Pro & $7.1\%$ & $1.1\%$ & $10.4\%$ \\
Gemini 2.5 Flash & $4.4\%$ & $1.7\%$ & $8.2\%$ \\
Gemini 2.5 Flash Lite & $2.0\%$ & $0.4\%$ & $4.5\%$ \\
Cosmos Reason 1 & $0.1\%$ & $0.5\%$ & $2.8\%$ \\
Cosmos Reason 2 & $0.9\%$ & $0.5\%$ & $1.2\%$ \\
\bottomrule
\end{tabular}

\caption{Interim Physion-Eval~\cite{physioneval2026} results. Panel~A reports three rates on the first $601$ release clips of a fixed randomized stream
($543$ with a reported glitch, $58$ without): the evidence-conditioned flag rate, the
anomaly-phrased finding rate, and, for reported-glitch clips only, the matched-report clip
rate; these measures are defined in the preceding text. Panel~B transcribes the benchmark's published critic and human
results from its Figure~4: $J$ is its Youden statistic averaged over the five generators,
and the flag rate is the share of generated clips called physically unrealistic, pooled
over the release's two viewpoints. The two panels use different samples and system
configurations and are juxtaposed descriptively, not compared under one protocol. A dash
in the matched-report row marks a rate its protocol does not define on unreported clips; $J$
values are absent from Panel~A because the required matched real-world reference videos
are withheld from the public release.}
\label{tab:physion}
\end{table}


\subsection{Recall by clip property}\label{sec:app-robust}
\Cref{tab:robust} groups the $304$ core flaws (run of \cref{tab:headline}) by severity, duration,
flaw count, and difficulty stratum. Subgroup recall stays within about ten points of the overall
rate everywhere, so neither severity nor stratum predicts whether a defect is found.

% GENERATED by analysis/emit_tex.py — do not edit by hand.
% Every number is recomputed from per-item rows; see the analysis scripts.

\begin{table}[t]
\centering\small
\setlength{\tabcolsep}{6pt}\renewcommand{\arraystretch}{1.12}
\begin{tabular}{@{}p{5.0cm} c c c@{}}
\toprule
\textbf{Group} & \textbf{Annotated} & \textbf{Found} & \textbf{Recall} \\
\midrule
\multicolumn{4}{@{}l}{\textit{Annotator severity}} \\
\quad severity 1 & 3 & 3 & 100\% \\
\quad severity 2 & 28 & 22 & 79\% \\
\quad severity 3 & 146 & 107 & 73\% \\
\quad severity 4 & 99 & 78 & 79\% \\
\quad severity 5 & 28 & 18 & 64\% \\
\addlinespace
\multicolumn{4}{@{}l}{\textit{Clip duration}} \\
\quad 6 s or less & 76 & 60 & 79\% \\
\quad 6-8 s & 46 & 29 & 63\% \\
\quad over 8 s & 182 & 139 & 76\% \\
\addlinespace
\multicolumn{4}{@{}l}{\textit{Number of annotated flaws per clip}} \\
\quad 1 & 63 & 49 & 78\% \\
\quad 2 & 92 & 65 & 71\% \\
\quad 3-4 & 113 & 87 & 77\% \\
\quad 5 or more & 36 & 27 & 75\% \\
\addlinespace
\multicolumn{4}{@{}l}{\textit{Difficulty stratum}} \\
\quad hard & 83 & 60 & 72\% \\
\quad not hard & 221 & 168 & 76\% \\
\bottomrule
\end{tabular}
\caption{Recall grouped by annotator severity, clip duration, number of annotated flaws in
the clip, and difficulty stratum, on the run of \cref{tab:headline}. Severity is the
annotator's own $1$--$5$ judgement of how badly the clip fails.}
\label{tab:robust}
\end{table}


\subsection{Refinement figures recomputed from the raw run}\label{sec:app-refine-verify}
\Cref{tab:refine-verify} recomputes every figure of \cref{sec:refine} from the per-clip records (all
reproduce exactly) and pins two definitions: the $8/502$ collateral figure and the differing baseline
and rewrite denominators.

% Independent recomputation check: every figure in the refinement tables,
% recomputed from the run's raw files (critic_full.jsonl, joint_rewritten.csv)
% pulled from S3, against the reported value. Recompute is stdlib-only, driven
% from per-clip rows, not from the summary.

\begin{table}[t]
\centering\small
\renewcommand{\arraystretch}{1.2}
\begin{tabularx}{\textwidth}{@{}L c c c@{}}
\toprule
\textbf{Quantity} & \textbf{Reported} & \textbf{Recomputed} & \\
\midrule
\multicolumn{4}{@{}l}{\textit{VeriPhy self-assessed metric, from the per-clip critic records ($589$ rows)}} \\
\quad acted (cited $\ge 1$ failure)              & $87/589$  & $87/589$  & \checkmark \\
\quad of acted, $p_{\text{plausible}}$ increased & $63/87$   & $63/87$   & \checkmark \\
\quad verdict implausible$\to$plausible          & $49/65$   & $49/65$   & \checkmark \\
\quad mean $\Delta p$ (acted)                    & $+0.524$  & $+0.524$  & \checkmark \\
\quad mean $p_{\text{plausible}}$ base$\to$rw    & $0.797\!\to\!0.862$ & $0.797\!\to\!0.862$ & \checkmark \\
\quad collateral (verdict flip plaus.$\to$implaus.) & $8/502$ & $8/502$   & \checkmark \\
\addlinespace
\multicolumn{4}{@{}l}{\textit{VideoPhy-2 AutoRater, rewrite side, per-clip scores ($589$ rows)}} \\
\quad joint \% (All / Hard)     & $23.1 / 3.4$  & $23.1 / 3.4$  & \checkmark \\
\quad mean SA (All / Hard)      & $3.13 / 2.76$ & $3.13 / 2.76$ & \checkmark \\
\quad mean PC (All / Hard)      & $3.58 / 3.32$ & $3.58 / 3.32$ & \checkmark \\
\quad SA\,$\ge 4$ \% (All / Hard) & $28.0 / 7.9$ & $28.0 / 7.9$ & \checkmark \\
\quad PC\,$\ge 4$ \% (All / Hard) & $54.5 / 36.2$ & $54.5 / 36.2$ & \checkmark \\
\bottomrule
\end{tabularx}
\caption{Every figure in \cref{tab:refine-selfmetric,tab:refine-autorater} recomputed
independently from the run's per-clip raw files, against the reported value.
All reproduce exactly. Two definitions are pinned by the recompute: \emph{collateral}
is a verdict flip (plausible$\to$implausible) on a non-flagged clip---$8/502$; scored
instead as \emph{any} decrease in $p_{\text{plausible}}$ it would be $31/502$, so the
stricter reading is the one reported. The AutoRater \emph{baseline} column of
\cref{tab:refine-autorater} is not in this check: the released bundle synced only the
rewrite-side per-clip scores and a $12$-clip baseline subset, so the $600$-clip baseline
means are corroborated by the run log rather than recomputed here. Baseline $n=600$ and
rewrite $n=589$ (the $11$ absent clips did not complete regeneration or re-scoring), so
the two columns of \cref{tab:refine-autorater} rest on slightly different denominators.}
\label{tab:refine-verify}
\end{table}


% =====================================================================
\renewcommand{\bibfont}{\footnotesize}
